% ============================================================================
\section{Pembahasan}
\label{sec:pembahasan}
% ============================================================================

\catatan{Subbagian A--C menunggu hasil. Subbagian D dan E sudah dapat ditulis dari rancangan, dan akan disesuaikan dengan hasil.}

\subsection{Jawaban atas pertanyaan penelitian}
\isian{Satu paragraf per RQ, berisi metrik utama dengan CI dan putusan hipotesisnya.}

\subsection{Analisis kesalahan LLM}
Taksonomi kesalahan ditetapkan sebelum pelabelan, agar label tidak menyesuaikan diri dengan keluaran yang sudah terlihat. Jenisnya adalah nilai salah, masa berlaku salah, kondisi atau cakupan salah, pengecualian terlewat, pencabutan terlewat, lapisan atau sifat salah, kutipan karangan, dan perubahan pada pengecoh. \isian{Frekuensi per jenis, gerbang yang menangkap setiap jenis, dan contoh kesalahan yang lolos semua gerbang.}

\subsection{Implikasi}
\isian{Implikasi bagi pengembang sistem payroll dan bagi formalisasi peraturan yang sering berubah.}

\subsection{Ancaman terhadap validitas}
\label{subsec:ancaman}

\textbf{Validitas internal.} \en{Ground truth} dan \en{engine} disusun tim yang sama, sehingga kesalahan yang sama dapat muncul di keduanya. Mitigasinya adalah acuan pihak ketiga, subset \en{holdout} contoh resmi, pengadjudikasi kedua, dan uji \en{engine} yang terpisah dari \en{knowledge base}.

\textbf{Bias pandangan belakang.} Representasi dirancang setelah perubahan 2023--2026 diketahui. Mitigasinya adalah perubahan tak terlihat yang ditulis orang lain, peristiwa prospektif, dan pelaporan riwayat \en{engine}.

\textbf{Validitas konstruk.} Kecocokan dengan acuan tidak sama dengan kesesuaian hukum. Penelitian ini tidak melibatkan pakar hukum pajak, kecuali bila praktisi pajak dapat dilibatkan untuk memeriksa sampel. Upaya manusia diukur dengan proksi, yaitu jumlah operasi sunting.

\textbf{Validitas eksternal.} Penelitian ini mencakup satu yurisdiksi, satu jenis pajak, dan satu rangkaian perubahan. Hasilnya belum tentu berlaku untuk pajak lain atau negara lain.

\textbf{Validitas kesimpulan.} Jumlah peristiwa nyata kecil, LLM tidak deterministik, dan versi model berubah dari waktu ke waktu. Karena itu hasil dilaporkan sebagai estimasi dengan selang kepercayaan eksak, dan seluruh masukan serta keluaran LLM diarsipkan.

\textbf{Kontaminasi data.} Regulasi 2023--2026 kemungkinan ada di data latih LLM. Mitigasinya adalah dokumen kontrafaktual, perbandingan hasil sebelum dan sesudah \en{training cutoff} model, dan peristiwa prospektif.

\subsection{Keterbatasan}
\label{subsec:keterbatasan}

Penelitian ini terbatas pada pegawai tetap dengan satu pemberi kerja. Pegawai tidak tetap, pensiunan, pemberi kerja ganda, dan fasilitas khusus Ibu Kota Nusantara tidak dicakup, dan nilai natura diterima sebagai masukan. Jenis perubahan yang tidak terjadi pada 2023--2026, seperti perubahan PTKP, hanya terwakili lewat data pengembangan dan perubahan tak terlihat. \isian{Keterbatasan lain dari hasil, misalnya jenis perubahan yang selalu membuat sistem \en{abstain}, dan peran independen yang tidak terisi.}
